\documentclass[11pt]{article}
\usepackage[letterpaper,margin=1in]{geometry}
\usepackage[T1]{fontenc}
\usepackage{lmodern,microtype}
\usepackage{amsmath,amssymb,amsthm,mathtools}
\usepackage{xcolor,hyperref,fancyhdr}
\hypersetup{colorlinks=true,linkcolor=blue!45!black,
 citecolor=blue!45!black,urlcolor=blue!45!black,
 pdftitle={Fixed-sample zero-free bounds for bounded-batch discovery},
 pdfauthor={Christopher D. Long}}
\pagestyle{fancy}\fancyhf{}
\fancyhead[L]{\small Fixed-sample bounds for batch discovery}
\fancyhead[R]{\small Research draft}\fancyfoot[C]{\thepage}
\setlength{\headheight}{14pt}\setlength{\emergencystretch}{2em}
\numberwithin{equation}{section}
\newtheorem{theorem}{Theorem}[section]
\newtheorem{lemma}[theorem]{Lemma}
\newtheorem{proposition}[theorem]{Proposition}
\newtheorem{corollary}[theorem]{Corollary}
\theoremstyle{remark}\newtheorem{remark}[theorem]{Remark}
\newcommand{\E}{\mathbb E}\newcommand{\Pp}{\mathbb P}
\newcommand{\R}{\mathbb R}\newcommand{\C}{\mathbb C}
\newcommand{\Var}{\operatorname{Var}}\newcommand{\ind}{\mathbf1}
\newcommand{\dK}{d_{\mathrm K}}\newcommand{\calZ}{\mathcal Z}
\newcommand{\Normal}{\mathcal N(0,1)}
\title{Fixed-sample zero-free bounds for bounded-batch discovery}
\author{Christopher D. Long}
\date{November 26, 2026\\[3pt]\small Research draft}
\begin{document}
\maketitle
\begin{abstract}
For discovery after exactly $n$ independent random batches, we prove that
the moment generating function is zero-free on $|z|<1/(4W)$ when each
batch has total species weight at most $W$. The proof replaces the
independent batch clocks of the Poissonized model by cliques of mutually
exclusive choices. A harmonic normal-approximation theorem of
Michelen and Sahasrabudhe then gives a uniform Kolmogorov bound
$CW\log(n+2)/s_n$, where $s_n^2$ is the actual fixed-sample variance.
The statements allow a countable species space, arbitrary overlap,
independent batches with different laws, and signed weights under an
absolute-weight bound. For iid nonnegative-weight models we also prove
uniform operator bounds for the canonical Hoeffding kernels. The optimal
order follows when $s_n^2\ge\eta nW^2$, with a constant depending on
$\eta>0$. A
two-community construction with batches of size two rules out any
universal zero-free sector about the entire positive real axis for the
probability generating function. This obstruction does not disprove the
conjectured optimal rate $CW/s_n$. Removing the logarithm remains open
here; a precise sufficient actual-variance moment estimate is identified.
\end{abstract}

\section{Results and scope}\label{sec:results}
Let $\mathcal S$ be countable, let $A_1,\ldots,A_n$ be independent random
finite subsets of $\mathcal S$, and let $h_i\in\R$ be deterministic weights.
The batch laws need not agree. Assume $n\ge1$, $W>0$, and
\begin{equation}\label{eq:budget}
 \sum_{i\in A_k}|h_i|\le W\quad\hbox{almost surely},\qquad 1\le k\le n.
\end{equation}
Define
\[
 D_{h,n}=\sum_{i\in\bigcup_{k=1}^n A_k}h_i,
 \qquad s_{h,n}^2=\Var D_{h,n},
 \qquad F_{h,n}(z)=\E e^{zD_{h,n}}.
\]
In particular, $|D_{h,n}|\le nW$, so $F_{h,n}$ is entire. For unweighted
discovery with $|A_k|\le B$, take $h_i=1$ and $W=B$. Write $d_{\mathrm K}$
for Kolmogorov distance and $Z\sim\Normal$.

\begin{theorem}[Fixed-sample zero-free disk]\label{thm:disk}
Under \eqref{eq:budget},
\begin{equation}\label{eq:disk}
 F_{h,n}(z)\ne0\qquad\left(|z|<\frac1{4W}\right).
\end{equation}
The disk is uniform in $n$, the batch laws, the number of species, and the
overlap pattern.
\end{theorem}

\begin{theorem}[A uniform rate with a logarithmic factor]\label{thm:lograte}
There is an absolute constant $C<\infty$ such that, whenever $s_{h,n}>0$,
\begin{equation}\label{eq:lograte}
 \dK\left(\frac{D_{h,n}-\E D_{h,n}}{s_{h,n}},Z\right)
 \le \min\left\{1,\frac{CW\log(n+2)}{s_{h,n}}\right\}.
\end{equation}
Consequently, for a sequence of such models, possibly with changing laws,
weights, and sample sizes, $s_{h,n}/[W\log(n+2)]\to\infty$ implies
asymptotic normality.
\end{theorem}

The normal-approximation input is the harmonic-function method of
Michelen--Sahasrabudhe \cite[Lemma~4.1 and Corollary~9.3]{MS}.
Section~\ref{sec:harmonic} supplies all model-specific hypotheses, including
a bounded-support argument that avoids a lattice assumption. A direct
application of their polynomial disk theorem would instead give
$CB\log(nB)/s_n$ in the unweighted case. Scaling by $W$ before applying
the harmonic argument gives \eqref{eq:lograte}.

The companion paper \cite{LongOptimal} proves the optimal-order
Poissonized bound $150W/\sqrt{V_h}$ for nonnegative weights. The present
fixed-sample result uses its own categorical representation and the
actual variance $s_{h,n}^2$; it does not condition a scalar Poissonized
normal approximation. The iid mean and residual-variance corrections in
\cite[Section~7]{LongQuantitative} are compatible with this result, but
do not by themselves prove \eqref{eq:lograte} or remove its logarithm.

\subsection*{The optimal target remains open}
Even for iid nonnegative-weight batches, we have not proved the uniform
bound
\begin{equation}\label{eq:target}
 \dK\left(\frac{D_{h,n}-\E D_{h,n}}{s_{h,n}},Z\right)
 \le\min\{1,CW/s_{h,n}\}.
\end{equation}
Section~\ref{sec:criterion} gives an explicit sufficient analytic estimate
for this target. Section~\ref{sec:hoeffding} proves an additional
operator bound, but does not assert that separate bounds on individual
Hoeffding orders control all their mixed cumulants. Section~\ref{sec:sector}
disproves a full-sector strengthening of \eqref{eq:disk}; it supplies no
counterexample to \eqref{eq:target}.

\section{An independence polynomial with categorical cliques}
\label{sec:cliques}
For a finite graph $G$ with activities $(\lambda_v)$, set
\[
 \calZ_G=\sum_{I\subseteq V(G):\ I\text{ independent}}
                  \prod_{v\in I}\lambda_v,
 \qquad \calZ_\varnothing=1.
\]
Every induced subgraph retains the original activities.

\begin{lemma}[Complex activities opposite disjoint cliques]\label{lem:cliques}
Suppose $V(G)=L\sqcup R$, $L$ is independent, and the induced graph on $R$
is a disjoint union of cliques. Give $i\in L$ activity $u_i\in\C$ and
$j\in R$ activity $w_j\ge0$. Suppose $0\le r_i<1$ satisfy
\begin{equation}\label{eq:arcsin}
 |u_i|\le r_i,\qquad
 \alpha_j:=\sum_{i\in N_G(j)\cap L}\arcsin r_i<\frac\pi4
 \quad(j\in R).
\end{equation}
Then $\calZ_G\ne0$. For every induced subgraph $H$ and vertex $v$ in $H$,
the ratio
\[
 R_v(H)=\lambda_v\frac{\calZ_{H-v-N_H(v)}}{\calZ_{H-v}}
\]
is defined and obeys
\begin{equation}\label{eq:ratios}
 |R_i(H)|\le r_i\quad(i\in L),\qquad
 \operatorname{Re}R_j(H)\ge0\quad(j\in R).
\end{equation}
\end{lemma}
\begin{proof}
We prove all assertions simultaneously by induction on the number of
vertices. The class of graphs in the statement is closed under taking
induced subgraphs. Assume the assertions for proper induced subgraphs
of a nonempty graph, again called $G$. In particular all proper-subgraph
partition functions used below are nonzero.

Fix $i\in L$ and list its neighbors as $j_1,\ldots,j_d$. The deletion
identity $\calZ_H=\calZ_{H-v}(1+R_v(H))$, successively applied to proper
subgraphs, gives
\[
 R_i(G)=u_i\prod_{\ell=1}^d(1+R_{j_\ell}(G-i-\{j_1,\ldots,j_{\ell-1}\}))^{-1}.
\]
Every factor in the denominator has modulus at least one by induction.
This proves the left inequality in \eqref{eq:ratios}.

Next fix $j\in R$ and let $C$ be its clique in $G[R]$. Put $H=G-C$.
For $k\in C$, define
\[
 q_k=\frac{\calZ_{H-(N_G(k)\cap L)}}{\calZ_H}.
\]
Delete these left neighbors one at a time within $H$. The resulting
telescoping expression for $q_k$ is a product of factors $(1+R_i)^{-1}$
in proper induced subgraphs. The elementary disk estimate
$|\arg(1+\zeta)|\le\arcsin r$ for $|\zeta|\le r<1$ yields
\[
 |\arg q_k|\le\alpha_k<\pi/4.
\]
There are no edges between distinct right cliques. Partitioning by the
at-most-one occupied vertex in $C$ therefore gives the exact identity
\begin{equation}\label{eq:clique-ratio}
 R_j(G)=\frac{w_jq_j}{1+\sum_{k\in C\setminus\{j\}}w_kq_k}.
\end{equation}
The denominator is in the open cone $|\arg z|<\pi/4$, which is convex
and contains the positive real axis. If $w_j>0$, the numerator is in the
same cone, so the ratio has argument strictly between $-\pi/2$ and
$\pi/2$. If $w_j=0$, the ratio is zero. This proves the right inequality
in \eqref{eq:ratios}.

For a left vertex, $|R_v(G)|<1$; for a right vertex,
$\operatorname{Re}(1+R_v(G))\ge1$. Thus $1+R_v(G)\ne0$ in either case.
The deletion identity proves $\calZ_G\ne0$, completing the induction.
\end{proof}

\begin{remark}
The angle budget is $\pi/4$, rather than the $\pi/2$ budget for the
bipartite lemma of \cite{LongOptimal}. In \eqref{eq:clique-ratio} both
the numerator and the denominator can rotate. Controlling each in the
$\pi/4$ cone keeps their quotient in the right half-plane. No degree
bound is placed on a species vertex.
\end{remark}

\section{From categorical choices to fixed-sample discovery}
\label{sec:representation}
First restrict to a finite species set $J$, intersecting each batch with
$J$. Each draw then has a finite law $\mu_k$ on subsets $C\subseteq J$;
empty choices are retained. For a positive parameter $M$, construct a
graph with left vertices $J$ and one right clique for each draw $k$.
The vertices in that clique are its supported options $(k,C)$, with
activities $M\mu_k(C)$. Join $(k,C)$ to each species in $C$. Assign
activity $u_i$ to left vertex $i$. Summing first over the left set gives
\begin{equation}\label{eq:finiteM}
 \calZ_M(\boldsymbol u)
 =\sum_{S\subseteq J}\left(\prod_{i\in S}u_i\right)
    \prod_{k=1}^n\left[1+M\Pp(A_k\cap S=\varnothing)\right].
\end{equation}
Consequently,
\begin{align}
 \lim_{M\to\infty}\frac{\calZ_M(\boldsymbol u)}{(1+M)^n}
 &=\sum_{S\subseteq J}\left(\prod_{i\in S}u_i\right)
           \prod_{k=1}^n\Pp(A_k\cap S=\varnothing)\notag\\
 &=\E\prod_{i\in J}(1+u_iY_i)=:Z_J(\boldsymbol u),
 \label{eq:absence}
\end{align}
where $Y_i$ indicates that species $i$ is absent after all $n$ draws.
Both the coefficient convergence and the convergence on compact sets
are immediate from the finite sum. Independence of the draws is exactly
what justifies the product of avoidance probabilities.

\begin{proof}[Proof of Theorem~\ref{thm:disk}]
For $|z|<1/(4W)$ put
\[
 u_i(z)=e^{-zh_i}-1,\qquad r_i=e^{|z||h_i|}-1.
\]
Species that never occur may be discarded. Every other species has
$|h_i|\le W$. For $0\le a\le1/4$,
$e^a-1\le2a$ and $\arcsin(e^a-1)\le(\pi/2)(e^a-1)\le\pi a$.
Thus $|u_i(z)|\le r_i<1$ and, for each supported option $C$,
\[
 \sum_{i\in C}\arcsin r_i
 \le\pi|z|\sum_{i\in C}|h_i|
 \le\pi|z|W<\pi/4.
\]
Lemma~\ref{lem:cliques} proves that
$\calZ_M(\boldsymbol u(z))/(1+M)^n$ is nonzero throughout this disk,
for every $M>0$. Its locally uniform limit in \eqref{eq:absence} is
one at $z=0$. Hurwitz's theorem shows that this limit is nonzero throughout
the disk. On the finite species set,
\[
 \E e^{zD_{h,n}^{(J)}}
   =e^{z\sum_{i\in J}h_i}\,Z_J((e^{-zh_i}-1)_{i\in J}),
\]
which proves the assertion there.

Take an increasing finite exhaustion $J\uparrow\mathcal S$.
The truncated discoveries tend almost surely to $D_{h,n}$ and are
bounded in absolute value by $nW$. Their MGFs converge locally uniformly
on $\C$, for example by dominated convergence applied to the supremum
on any compact $z$-set. A second application of Hurwitz, again using
the value one at zero, proves \eqref{eq:disk} on the countable space.
\end{proof}

\section{A harmonic bound using the support scale}\label{sec:harmonic}
We record precisely the external analytic inputs. For $R>1$, set
\[
 S_R(\delta)=\{q:R^{-1}\le|q|\le R,\ |\arg q|\le\delta\},
\]
and let $S_R^*(\delta)$ be the two radial ends $|q|\in\{R^{-1},R\}$.
A real harmonic function $u$ is symmetric if $u(\bar q)=u(q)$, and
weakly positive if $u(|q|)\ge u(q)$. It is $b$-decreasing on a region if
\[
 u(\rho e^{i\theta_1})-u(\rho e^{i\theta_2})\ge-b
 \quad(0\le\theta_1\le\theta_2)
\]
whenever the two points belong to that region.

Michelen--Sahasrabudhe \cite[Lemma~4.1]{MS} prove that a symmetric,
weakly positive harmonic $u$ on a neighborhood of $S_R(\delta)$ is
$b$-decreasing on $S_r(\delta/2)$ if $R>r>1$ and
\begin{equation}\label{eq:MS-boundary}
 (r/R)^{1/\delta}\max_{S_R^*(\delta)}|u|\le3b/8.
\end{equation}
Their Corollary~9.3 states that, for a random variable with standard
deviation $\sigma$, a harmonic $b$-decreasing logarithmic potential
$u(q)=\log|\E q^Y|$ on $B(1,\epsilon)$ gives Kolmogorov error
$O((\epsilon\sigma)^{-1})$ if
\begin{equation}\label{eq:MS-cumulants}
 \sum_{j\ge2}\frac{|\kappa_j(Y)|}{j!}
                   (\epsilon/32)^j>b.
\end{equation}
These analytic statements apply to real-valued variables with exponential
moments, using a logarithm branch near one; see also \cite[Section~12]{MS}.

\begin{lemma}[Bounded support and a fixed MGF disk]\label{lem:harmonic}
Let $Y\in[0,N]$ almost surely, where $N\ge1$, and let
$\sigma^2=\Var Y>0$. If $\E e^{zY}\ne0$ for $|z|<1/4$, then
\begin{equation}\label{eq:abstract-log}
 \dK((Y-\E Y)/\sigma,Z)
 \le\min\{1,C\log(N+2)/\sigma\}
\end{equation}
for an absolute constant $C$.
\end{lemma}
\begin{proof}
Let $K(z)=\log\E e^{zY}$ be the analytic branch with $K(0)=0$.
On $|z|<1/4$, the modulus inequality gives
$\operatorname{Re}K(z)\le N\max(\operatorname{Re}z,0)<N/4$.
Apply Harnack's inequality to the positive harmonic function
$H(z)=N/4-\operatorname{Re}K(z)$, for which $H(0)=N/4$.
On $|z|\le1/8$, it gives $H(z)\le3N/4$, and hence
\begin{equation}\label{eq:harnack}
 |\operatorname{Re}K(z)|\le N/2\qquad(|z|\le1/8).
\end{equation}
Use the principal logarithm and put
$f(q)=\E q^Y=\exp K(\log q)$ and $u(q)=\log|f(q)|$ near one.
This $u$ is harmonic, symmetric, and weakly positive; the last property
follows from $|\E q^Y|\le\E|q|^Y$.

Choose
\[
 R=e^{1/16},\quad r=e^{1/32},\quad
 \delta=\frac1{64\log(2N+2)},\quad \epsilon=\delta/4.
\]
The whole sector $S_R(\delta)$ satisfies $|\log q|<1/8$, with room to
spare, so \eqref{eq:harnack} bounds $|u|$ there by $N/2$. Moreover,
\[
 (r/R)^{1/\delta}\max_{S_R^*(\delta)}|u|
 \le\frac{N}{2(2N+2)^2}\le\frac1{32}<\frac38.
\]
The input \eqref{eq:MS-boundary} with $b=1$ proves that $u$ is
1-decreasing on $S_r(\delta/2)$. The disk $B(1,\epsilon)$ lies inside
that sector: its angular deviation is at most $2\epsilon=\delta/2$,
and its radial range lies between $e^{-1/32}$ and $e^{1/32}$.

If $\sigma>64/\epsilon$, the $j=2$ term alone gives
\[
 \sum_{j\ge2}\frac{|\kappa_j(Y)|}{j!}(\epsilon/32)^j
 \ge\frac{\sigma^2\epsilon^2}{2048}>1.
\]
Apply \eqref{eq:MS-cumulants} to obtain $C/(\epsilon\sigma)$.
If $\sigma\le64/\epsilon$, enlarge the absolute constant and use
$d_{\mathrm K}\le1$. Since $\epsilon^{-1}=256\log(2N+2)$, this proves
\eqref{eq:abstract-log}.
\end{proof}

\begin{proof}[Proof of Theorem~\ref{thm:lograte}]
For nonnegative weights, apply Lemma~\ref{lem:harmonic} to
$Y=D_{h,n}/W\in[0,n]$. Theorem~\ref{thm:disk} gives the required disk,
and $\sigma=s_{h,n}/W$. For signed weights use
$Y=(D_{h,n}+nW)/W\in[0,2n]$. Translation multiplies the MGF by a nonzero
exponential and does not change the standardized distribution. Since
$\log(2n+2)\le2\log(n+2)$, the same conclusion holds after adjusting $C$.
\end{proof}

\section{Operator bounds for the iid Hoeffding kernels}
\label{sec:hoeffding}
This section assumes iid draws and nonnegative weights. It supplies
additional variance-sensitive structure, without claiming the optimal rate.
Let $p_i=\Pp(i\in A_1)$, $q_i=1-p_i$, and
$\xi_i(A)=\ind_{\{i\in A\}}-p_i$. Initially take finitely many species.
Let $\mathcal H=L_0^2(\mu)$ be the real Hilbert space of centered
square-integrable functions of one batch. For $1\le k\le n$, define the
canonical symmetric kernel in $\mathcal H^{\otimes k}$ by
\begin{equation}\label{eq:kernel}
 H_k=(-1)^{k+1}\sum_i h_iq_i^{n-k}\xi_i^{\otimes k}.
\end{equation}
As usual a zero-th power is one, including at $q_i=0$.
Expanding the species-absence products gives the orthogonal decomposition
\begin{equation}\label{eq:Hoeffding}
 D_{h,n}-\E D_{h,n}
 =\sum_{k=1}^n\sum_{\substack{J\subseteq[n]\\|J|=k}}H_k(A_J),
 \qquad
 s_{h,n}^2=\sum_{k=1}^n\binom nk\|H_k\|_2^2.
\end{equation}
Orthogonality follows by conditioning on a coordinate in the symmetric
difference of two subsets. The first kernel is a centered function whose
range has length at most $W$.

\begin{proposition}[Every nontrivial flattening is bounded]\label{prop:flatten}
For $2\le k\le n$ and $1\le r<k$, regard $H_k$ as an operator from
$\mathcal H^{\otimes r}$ to $\mathcal H^{\otimes(k-r)}$. Then
\begin{equation}\label{eq:flatten}
 \|H_k\|_{r\to k-r}
 \le W\sup_i p_i^{k/2}q_i^{n-k},\qquad
 \sqrt{\binom nk}\,\|H_k\|_{r\to k-r}\le W.
\end{equation}
These inequalities and \eqref{eq:Hoeffding} extend to the countable species
space by $L^2$ limits.
\end{proposition}
\begin{proof}
Species with $p_i=0$ contribute nothing and are discarded. For $r\ge1$
and $g\in\mathcal H^{\otimes r}$, define
\[
 (T_rg)_i=\sqrt{h_i}\,p_i^{-r/2}
                    \langle g,\xi_i^{\otimes r}\rangle.
\]
Canonicality in every coordinate gives
\[
 \langle g,\xi_i^{\otimes r}\rangle
   =\E\left[g(A_1,\ldots,A_r)
                    \prod_{a=1}^r\ind_{\{i\in A_a\}}\right].
\]
Cauchy--Schwarz, followed by independence of these $r$ batches, implies
\begin{align*}
 \|T_rg\|_{\ell^2}^2
 &\le\sum_i h_i\E\left[g^2
                    \prod_{a=1}^r\ind_{\{i\in A_a\}}\right]\\
 &=\E\left[g^2\sum_{i\in\bigcap_{a=1}^r A_a}h_i\right]
 \le W\|g\|_2^2.
\end{align*}
Thus $\|T_r\|\le\sqrt W$. The flattening of \eqref{eq:kernel} factors,
up to its sign, as
\[
 T_{k-r}^*\operatorname{diag}(p_i^{k/2}q_i^{n-k})T_r.
\]
This proves the first inequality. For every $p\in[0,1]$,
\[
 \binom nk p^k(1-p)^{2(n-k)}
 \le\binom nk p^k(1-p)^{n-k}\le1,
\]
which proves the second.

For a countable space, boundedness by $nW$ gives $L^2$ convergence of
truncated discoveries. Orthogonal Hoeffding projections are continuous
in $L^2$, so each kernel converges in Hilbert--Schmidt norm, hence also
in every flattening operator norm. Pass to the limit in the inequalities.
\end{proof}

\begin{corollary}[Summable operator bounds across orders]\label{cor:summable}
For $2\le k\le n$, put
$a_k=\sqrt{\binom nk}\max_{1\le r<k}\|H_k\|_{r\to k-r}$.
Then
\[
 a_k\le W(2/3)^{k/2}+2W2^{-n/2}.
\]
In particular, for every fixed $a\ge0$ there is a finite constant $C_a$
such that $\sum_{k=2}^n k^a a_k\le C_aW$, uniformly in $n$ and the law.
\end{corollary}
\begin{proof}
Split the species at $p_i=1/2$. On the low-probability part,
with $q=1-p\ge1/2$,
\[
 \binom nk p^kq^{2(n-k)}
 =(1+q)^{-k}\Pp\{\operatorname{Bin}(n,1-q^2)=k\}
 \le(2/3)^k.
\]
Proposition~\ref{prop:flatten} therefore gives the first term. The
high-probability species have total weight at most $2W$, because
$\sum_i h_ip_i=\E\sum_{i\in A_1}h_i\le W$.
The rank-one kernel $\xi_i^{\otimes k}$ has every flattening norm
$\|\xi_i\|_2^k=(p_iq_i)^{k/2}$. Thus the triangle inequality gives
\begin{align*}
 \sqrt{\binom nk}\,\|H_k^{\rm high}\|_{r\to k-r}
 &\le\sum_{p_i>1/2}h_iq_i^{n/2}
       \sqrt{\binom nk p_i^kq_i^{n-k}}\\
 &\le2W2^{-n/2}.
\end{align*}
Combine the two parts. The series $\sum_{k\ge2}k^a(2/3)^{k/2}$
converges, and $2^{-n/2}\sum_{k=2}^n k^a$ is uniformly bounded.
\end{proof}

An operator norm controls contractions of a fixed kernel. It is not the
Hilbert--Schmidt norm in the variance identity, and is not a bound for
all mixed contractions across different values of $k$. Although
Corollary~\ref{cor:summable} sums these operator bounds with any fixed
polynomial weight, a variance-sensitive bound for the mixed contractions
themselves remains to be proved. No such estimate
is assumed in Theorems~\ref{thm:disk}--\ref{thm:lograte}.

\section{A full zero-free sector is impossible}\label{sec:sector}
The sector theorem of \cite[Theorem~1.4]{MS} removes the logarithm under
a stronger hypothesis than a disk. That hypothesis fails uniformly here.

\begin{theorem}[Two-community sector obstruction]\label{thm:sector}
There is no $c>0$ such that the probability generating function of every
iid fixed-sample discovery count with batch bound $B$ is zero-free for
$w\ne0$ with $|\arg w|<c/B$. This fails already with $B=2$.
\end{theorem}
\begin{proof}
Fix $M\ge3$ and two disjoint communities of $M$ species each. One draw
chooses a community with probability $1/2$, then a uniformly random
unordered pair within it. Let $G_{M,n}(w)=\E w^{D_n}$. We prove the
locally uniform polynomial limit
\begin{equation}\label{eq:root-limit}
 2^{2n}G_{M,n}(2^{-n/M}z)\longrightarrow z^{2M}+2z^M
 \qquad(n\to\infty,\ M\text{ fixed}).
\end{equation}

For a subset with $a$ species from the first community and $b$ from the
second, the probability that one batch stays inside it is
\[
 q_{a,b}=\frac{\binom a2+\binom b2}{2\binom M2}.
\]
Writing $x=a/M$, $y=b/M$, one has
\begin{equation}\label{eq:capacity}
 q_{a,b}\le\frac{x^2+y^2}{2}\le2^{x+y-2}.
\end{equation}
For the first inequality use
$a(a-1)/(M(M-1))\le(a/M)^2$. For the second fix $s=x+y$.
If $0\le s\le1$, the largest possible $x^2+y^2$ is $s^2$, and
$s^2\le2^{s-1}$: on $(0,1]$, the logarithm of their ratio has positive
derivative $2/s-\log2$ and is zero at one. If $1\le s\le2$, put
$u=s-1$. The largest possible sum is $1+u^2\le2^u$, because
$\log(1+u^2)$ is convex on $[0,1]$ and lies below its endpoint chord.
Equality in \eqref{eq:capacity} is possible only at
\[
 (a,b)=(M,0),(0,M),(M,M).
\]

The event that the observed union equals a given subset $T$ implies
that every draw stays in $T$. Summing over all subsets of size $d$ and
using strictness in \eqref{eq:capacity} shows
\[
 2^{2n-nd/M}\Pp(D_n=d)\longrightarrow0
 \quad(d\notin\{M,2M\}).
\]
There are only finitely many subsets for fixed $M$. Every species has
incidence probability $1/M$, so a union bound gives
$\Pp(D_n<2M)\le2M(1-1/M)^n\to0$. Finally, all draws choosing one fixed
community has probability $2^{-n}$; conditional on this event that
community is fully covered with probability tending to one. The two
such events give $\Pp(D_n=M)=2^{1-n}(1+o(1))$. Any other union of size
$M$ has capacity strictly below $1/2$ and contributes an exponentially
smaller amount. This proves \eqref{eq:root-limit} coefficientwise,
hence locally uniformly.

The nonzero roots of the limit are simple. In particular,
$\zeta_M=2^{1/M}e^{i\pi/M}$ is a simple root. Rouch\'e's theorem on a
small circle about it yields a root $w_{M,n}$ of $G_{M,n}$ with
\begin{equation}\label{eq:smallangle}
 w_{M,n}\sim2^{-(n-1)/M}e^{i\pi/M}.
\end{equation}
Given $c>0$, choose $M>2\pi/c$ and then $n$ large. The root has
$|\arg w_{M,n}|<c/2=c/B$, as required.
\end{proof}

\begin{remark}[The obstruction does not refute the optimal rate]
For fixed $M$, put $q=1-1/M$. Individual absences give
$\E(2M-D_n)=2Mq^n$. Every pair of distinct species has joint absence
probability $\rho^n$ with $\rho<q$. Since $M$ is fixed, higher factorial
moments of the number absent are $o(q^n)$. Consequently
\[
 \Var D_n\sim2Mq^n,\qquad \kappa_3(D_n)\sim-2Mq^n.
\]
The untilted variance tends to zero, making \eqref{eq:target} vacuous.
The corresponding MGF roots satisfy
$\log w_{M,n}=-(n-1)(\log2)/M+i\pi/M+o(1)$ and are far from zero.
Thus they are consistent with Theorem~\ref{thm:disk}. Exponentially
tilting this count does not produce another iid batch model by definition;
such a tilt cannot be used without proof to construct a rate counterexample.
\end{remark}

\section{What would remove the logarithm?}\label{sec:criterion}
\begin{proposition}[A sufficient actual-variance estimate]\label{prop:criterion}
Consider a class of models covered by Theorem~\ref{thm:disk}.
Suppose there are absolute $A,a>0$ such that every model in the class
satisfies, with $D=D_{h,n}$ and $s^2=\Var D>0$,
\begin{equation}\label{eq:missing-mgf}
 \log\E e^{\theta(D-\E D)}\le A s^2\theta^2
 \qquad (\theta\in\R,\ |\theta|\le a/W).
\end{equation}
Then the optimal-order conclusion \eqref{eq:target} holds throughout
that class, with a constant depending only on $A,a$.
\end{proposition}
\begin{proof}
Put $c_0=\min(a/2,1/8)$, $r=c_0/W$, and let $K$ be the analytic
logarithm of the centered MGF on the zero-free disk. On $|z|\le r$,
the modulus inequality and \eqref{eq:missing-mgf} give
$\operatorname{Re}K(z)\le A s^2r^2$. Since $K(0)=0$, the positive
harmonic function $A s^2r^2-\operatorname{Re}K(z)$ bounds the Fourier
coefficients on the radius-$r$ circle. Explicitly, for $k\ge2$,
\[
 \frac{|\kappa_k(D)|}{k!}\le2A s^2r^{2-k}.
\]
For $k\ge3$ these are variance-linear cumulant bounds with scale
$W/c_0$. Summing the Taylor remainder for $K(it/s)$ up to
$|t|\le sr/[16(A+1)]$ bounds it by $C_A|t|^3/(sr)$ and by $t^2/4$.
Thus the characteristic-function difference is at most
$C_A|t|^3e^{-t^2/4}/(sr)$ in that range. Esseen smoothing, as used in
\cite[Section~5]{LongOptimal}, gives $C_A/(sr)=C(A,a)W/s$.
Use the trivial bound one when this exceeds one.
\end{proof}

\begin{corollary}[Optimal order when variance is proportional to sample size]
\label{cor:nondegenerate}
Fix $\eta>0$. For the models of Theorem~\ref{thm:disk} satisfying
$s_{h,n}^2\ge\eta nW^2$, there is a constant $C_\eta$, independent of
the remaining parameters, such that
\[
 \dK\left(\frac{D_{h,n}-\E D_{h,n}}{s_{h,n}},Z\right)
 \le\min\{1,C_\eta W/s_{h,n}\}.
\]
\end{corollary}
\begin{proof}
With all other draws fixed, changing one batch changes a nonnegative
weighted discovery by at most $W$, or a signed discovery by at most
$2W$. The successive Doob differences therefore have conditional ranges
of length at most $W$ or $2W$, respectively. The conditional Hoeffding
lemma and iteration give, for every real $\theta$,
\[
 \log\E e^{\theta(D-\E D)}
 \le nW^2\theta^2/8\quad\hbox{or}\quad nW^2\theta^2/2.
\]
Under the variance assumption this is \eqref{eq:missing-mgf} with
$A=1/(8\eta)$ or $1/(2\eta)$ and, for example, $a=1$.
Apply Proposition~\ref{prop:criterion}.
\end{proof}

Inequality \eqref{eq:missing-mgf} is a sufficient condition, not a
proved input for the full fixed-sample class. The corollary covers a
specified regime and does not control the dependence on $\eta$ as it
tends to zero. The Poissonized real MGF bound controls its
Poissonized variance. Its signed residual after removing the sample-count
direction is a different random variable. Likewise, a bounded-difference
estimate with proxy $nW^2$ or a self-bounding estimate with proxy
$W\E D$ cannot simply replace $s^2$ in \eqref{eq:missing-mgf}.
Michelen--Sahasrabudhe \cite[Theorem~11.4]{MS} show that a disk hypothesis
alone can require a logarithm for general distributions; removing it here
requires additional batch structure.

\begin{proposition}[The target order cannot be improved uniformly]
\label{prop:lower}
There are iid models with batch bound $B$, $s_n/B\to\infty$, and
\[
 \dK((D_n-\E D_n)/s_n,Z)\ge\frac{B}{16s_n}.
\]
Thus $B/s_n$ is a necessary order for a uniform fixed-sample theorem.
\end{proposition}
\begin{proof}
Choose uniformly among $m$ disjoint batches, each of size $B$, and take
$n=m$. Then $D_n=BK_m$, where $K_m$ is the number of occupied boxes in
$m$ uniform singleton draws into $m$ boxes. Its elementary indicator
variance formula gives
\[
 \Var K_m=mq_m(1-q_m)+m(m-1)(r_m-q_m^2)
 \sim m(e^{-1}-2e^{-2}),
\]
where $q_m=(1-1/m)^m$ and $r_m=(1-2/m)^m$.
In particular $s_n/B\to\infty$. Chebyshev's inequality places at least
$3/4$ of the mass of $D_n$ in an interval of length $4s_n$. Such an
interval contains at most $4s_n/B+2$ points of the lattice $B\mathbb Z$.
The largest atom is therefore at least $3/[4(4s_n/B+2)]$.
Any jump of mass $p$ has Kolmogorov distance at least $p/2$ from a
continuous distribution. For $s_n\ge B$, the asserted bound follows.
\end{proof}

\section*{Verification and research status}
The repository includes deterministic finite calculations for the
categorical partition identity, induced-subgraph ratios, harmonic
parameters, and the two-community coefficient limit. They supplement
the proofs and do not certify nonvanishing on a continuum, infinite-space
limits, or universal normal-approximation inequalities. The separate
internal proof audit records the hypotheses and the unresolved optimal
step. No peer-review, formal-verification, or exhaustive priority claim
is made for this research draft.

\begin{thebibliography}{9}
\bibitem{MS}
M.~Michelen and J.~Sahasrabudhe,
\emph{Central limit theorems and the geometry of polynomials},
J. Eur. Math. Soc. \textbf{28} (2026), 2261--2305.
\href{https://doi.org/10.4171/JEMS/1530}{doi:10.4171/JEMS/1530}.
\bibitem{LongOptimal}
C.~D.~Long, \emph{Optimal Berry--Esseen bounds for bounded-batch discovery},
research draft, November 26, 2026.
\url{https://github.com/long-mathematics/bounded-batch-discovery/blob/main/papers/optimal_berry_esseen.tex}.
\bibitem{LongQuantitative}
C.~D.~Long, \emph{Further quantitative bounds for bounded-batch discovery},
research draft, November 26, 2026.
\url{https://github.com/long-mathematics/bounded-batch-discovery/blob/main/papers/quantitative_extensions.tex}.
\end{thebibliography}
\end{document}
